\begin{thebibliography}{99}

\bibitem{devfuria}
DEVFURIA.
\textit{Introdução ao algoritmo de ordenação Insertion Sort}.
Disponível em:
\url{http://devfuria.com.br/logica-de-programacao/introducao-ao-algoritmo-de-ordenacao-insertion-sort/}.
Acesso em: 14 set. 2026.

\bibitem{cplusplusrand}
CPLUSPLUS.COM.
\textit{rand -- C++ Reference}.
Disponível em:
\url{https://cplusplus.com/reference/cstdlib/rand/}.
Acesso em: 14 set. 2026.

\bibitem{guru99}
GURU99.
\textit{C++ Manipulação de arquivos: como abrir, gravar, ler e fechar arquivos em C++}.
Disponível em:
\url{https://www.guru99.com/pt/cpp-file-read-write-open.html}.
Acesso em: 14 set. 2026.

\bibitem{medirtempo}
NGHIATU.
\textit{8 maneiras de medir o tempo de execução em C / C++}.
Disponível em:
\url{https://ichi.nghiatu.com/pt/8-maneiras-de-medir-o-tempo-de-execucao-em-c-c-79591185633529}.
Acesso em: 14 set. 2026.

\bibitem{stackoverflowparametros}
STACK OVERFLOW EM PORTUGUÊS.
\textit{Como criar parâmetros opcionais em C++?}.
Disponível em:
\url{https://pt.stackoverflow.com/questions/514913/como-criar-par%c3%a2metros-opcionais-em-c}.
Acesso em: 27 set. 2026.

\bibitem{shellsort}
SCRIBD.
\textit{Método de Shell Sort Explicado}.
Disponível em:
\url{https://pt.scribd.com/document/945722840/Shell-Sort}.
Acesso em: 27 set. 2026.

\bibitem{bubblesort}
TRYBE.
\textit{Bubble Sort: o que é e como usar? Exemplos práticos!}.
Disponível em:
\url{https://blog.betrybe.com/tecnologia/bubble-sort-tudo-sobre/}.
Acesso em: 27 set. 2026.

\bibitem{selectionsort}
JOÃO ARTHUR BARBOSA MOURA.
\textit{Estruturas de Dados e Algoritmos}.
Disponível em:
\url{https://joaoarthurbm.github.io/eda/posts/selection-sort/}.
Acesso em: 27 set. 2026.

\bibitem{comparacaoalgoritmos}
GEEKSFORGEEKS.
\textit{Comparison among Bubble Sort, Selection Sort and Insertion Sort}.
Disponível em:
\url{https://www.geeksforgeeks.org/dsa/comparison-among-bubble-sort-selection-sort-and-insertion-sort/}.
Acesso em: 27 set. 2026.

\bibitem{selectionbubble}
EDUCBA.
\textit{Selection sort vs Bubble sort | Differences of Selection sort vs Bubble sort}.
Disponível em:
\url{https://www.educba.com/selection-sort-vs-bubble-sort/}.
Acesso em: 27 set. 2026.

\bibitem{stackoverflow}
STACK OVERFLOW EM PORTUGUÊS.
\textit{Como passar uma função como parâmetro em C?}.
Disponível em:
\url{https://pt.stackoverflow.com/questions/2983/como-passar-uma-fun%C3%A7%C3%A3o-como-par%C3%A2metro-em-c}.
Acesso em: 14 set. 2026.

\bibitem{visualgo}
VISUALGO.
\textit{Sorting}.
Disponível em:
\url{https://visualgo.net/en/sorting}.
Acesso em: 27 set. 2026.

\bibitem{programizsorting}
PROGRAMIZ.
\textit{Sorting Algorithms}.
Disponível em:
\url{https://www.programiz.com/dsa/sorting-algorithm}.
Acesso em: 27 set. 2026.

\bibitem{programizshell}
PROGRAMIZ.
\textit{Shell Sort}.
Disponível em:
\url{https://www.programiz.com/dsa/shell-sort}.
Acesso em: 27 set. 2026.

\bibitem{mergegiovani}
GIOVANI DA CRUZ.
\textit{Merge Sort: um método eficaz de ordenação}.
Disponível em:
\url{https://giovanidacruz.com.br/merge-sort-um-metodo-eficaz-de-ordenacao/}.
Acesso em: 24 out. 2026.

\bibitem{geeksmerge}
GEEKSFORGEEKS.
\textit{Merge Sort}.
Disponível em:
\url{https://www.geeksforgeeks.org/dsa/merge-sort/}.
Acesso em: 24 out. 2026.

\bibitem{coddymerge}
CODDY.
\textit{Merge Sort — Visualização de algoritmo de ordenação}.
Disponível em:
\url{https://coddy.tech/visualize/pt/sorting/merge-sort}.
Acesso em: 24 out. 2026.

\bibitem{geeksquick}
GEEKSFORGEEKS.
\textit{Quick Sort}.
Disponível em:
\url{https://www.geeksforgeeks.org/dsa/quick-sort-algorithm/}.
Acesso em: 9 nov. 2026.

\bibitem{imequick}
FEOFILOFF, Paulo.
\textit{Algoritmo Quicksort: análise da correção e do desempenho}.
Instituto de Matemática e Estatística, Universidade de São Paulo.
Disponível em:
\url{https://users.ime.usp.br/~pf/algoritmos/aulas/quick.html}.
Acesso em: 8 nov. 2026.

\bibitem{programizquick}
PROGRAMIZ.
\textit{QuickSort (With Code in Python/C++/Java/C)}.
Disponível em:
\url{https://www.programiz.com/dsa/quick-sort}.
Acesso em: 9 nov. 2026.

\bibitem{ia}
OPENAI.
\textit{ChatGPT}: ferramenta de inteligência artificial utilizada como
auxílio na elaboração e revisão dos scripts de geração de gráficos e tabelas e organização LaTex no OverLeaf na elaboração do relatório.
2026.


\end{thebibliography}